\section{Cayley-Hamilton Theorem}

The characteristic equation for eigenvalues $0 = \det|{A - \lambda I |}$ can be expanded as a polynomial in $\lambda$ 
\[ \phi(\lambda) = \det{|A - \lambda I|} = a_n\lambda^n + a_{n-1}\lambda^{n-1} +\hdots + a_1 \lambda + a_0, \quad a_n = (-1)^n.\]
And the characteristic equation can be written in the form $\phi(\lambda) = 0$.\\

If in any polynomial $f(\lambda) = \sum c_i \lambda^i$ we substitute the matrix $A$ for $\lambda$ and multiply the constant term
$c_0$ by $I$, we obtain a matrix denoted by $f(A)$.

\begin{theorem} (Cayley-Hamilton) Any square matrix satisfies its characteristic equation. \end{theorem}

\begin{proof}
Let $\phi(\lambda)$ be the characteristic determinant of $A$. We want to show that $\phi(A) = 0$.  Since the elements of $A - \lambda I$ are 
linear functions of $\lambda$, and the elements of its adjoint [the transpose of the matrix of signed minors of $A-\lambda I$ -- denoted by Adj\_$(A-\lambda I)$] are (n-1)-rowed  determinants,
they are polynomials in $\lambda$ of degree $\le n-1$. Let $C$ stand for the Adj-$(A-\lambda I)$.
If the element in the i-th row and j-th column of $C$ is $\sum_k c_{ijk}\lambda^k$, then 

\[ C = \sum_{k=0}^{n-1} C_k \lambda^k, \quad C_k = (c_{ijk})\quad (i,j = 1, 2, \hdots, n) \]
 
 Using the relationship between the adjoint and the determinant, $\phi(\lambda) = \det{|A - \lambda I|}$, we can write
 
 \[ (A - \lambda I) C = \phi(\lambda) I. \] This expression can be expanded as 
 \[ A \sum_{k = 0} ^{n-1} C_k \lambda^k  - \lambda \sum_{k=0}^{n-1} C_k \lambda^k = \sum_{k=0}^n a_k \lambda^k  I \]
 
 Equating the terms free of $\lambda$ and the coefficients of $\lambda, \lambda^1, ..., \lambda^{n-1}, \lambda^n$ we get
 
 \begin{eqnarray*}
 AC_0                      &=& a_0 I,\\
 AC_1 - C_0            &=& a_1 I,\\
 AC_2 - C_1            &=& a_2 I,\\
 \hdots                     &~& \hdots \\
 AC_{n-1} - C_{n-2} &=& a_{n-1} I,\\
-C_{n-1}                  &=& a_n I.
\end{eqnarray*}

Multiply these equations on the left by $I, A, A^2, ..., A^{n-1}, A^n$ respectively and add; we get

\[ 0 = a_0 I + a_1 A + a_2 A^2 + \hdots + a_{n-1}A^{n-1} + a_n A^n = \phi(A), \] and we see that the matrix $A$ satisfies its own
characteristic equation.  \qed
\end{proof}

\section{Sylvesters Criterion for Positive Definite Hermitian Matrices}

\begin{theorem} (Sylvester's Criterion)  An $n$-x-$n$ Hermitian matrix M is positive-definite if and only if all the leading principal minors are positive. 
\end{theorem}

\begin{proof}
Let $M$ be an n x n Hermitian positive definite matrix. Then $0 \le x^\dagger M x$ if $x$ is not the zero vector. Let $M_k$ be the $k$-th leading principal minor 
(the first $k$-rows and $k$-columns in the upper left hand corner of $M$). Select a non-zero vector $x$ as follows 

\[ x = 
\begin{bmatrix}
x_1 \\
\vdots \\
x_k\\
0\\
\vdots\\
0
\end{bmatrix}
= 
\begin{bmatrix}
\vec{x}\\
0\\
\vdots\\
0
\end{bmatrix}
\]

Then we have that $0 \le x^\dagger M x = \vec{x}^\dagger M_k \vec{x}$. Therefore, since $\vec{x} \ne \vec{0}$ is otherwise arbitrary, then all the eigenvalues of $M_k$ must be
positive and hence  $\det|M_k| > 0 $

To prove the reverse implication, we use induction. The general form of an $(n+1)$-x-$(n+1)$ Hermitian matrix is

\[ M_{n+1} = 
\begin{bmatrix}
M_n & \vec{v}\\
\vec{v}^\dagger & d
\end{bmatrix}
\]

Denote the $(n+1)$ dimensional vector $x$ as 

\[ x = 
\begin{bmatrix}
\vec{x} \\
x_{n+1}
\end{bmatrix}
\]

Then 

\[ x^\dagger M_{n+1} x = \vec{x}^\dagger M_n \vec{x} + x_{n+1} \vec{x}^\dagger \vec{v} + \overline{x}_{n+1} \vec{v}^\dagger \vec{x} + d |x_{n+1}|^2 \]

By completing the square, this last expression is equal to 
\[ (\vec{x}^\dagger + \vec{v}^\dagger M_n^{-1} \overline{x}_{n+1}) M_n (\vec{x} +  x_{n+1}M_n^{-1} \vec{v} ) - |x_{n+1}|^2  \vec{v}^\dagger M_n^{-1} \vec{v} + d|x_{n+1}|^2 \]
\[ = (\vec{x} + \vec{c})^\dagger M_n (\vec{x} + \vec{c}) + |x_{n+1}|^2 (d - \vec{v}^\dagger M_n^{-1} \vec{v}), \] where $\vec{c} = x_{n+1} M_n^{-1}\vec{v}$. $M_n^{-1}$ exists because
the eigenvalues of $M_n$ are all positive. The first term is positive by the induction hypothesis.  Using this useful determinant expansion

\[ \det{ \begin{bmatrix}
A & B  \\
C & D \\
\end{bmatrix}} = \det{|A|} \det{ | D - C A^{-1}B | } \]

We have that 
\[ \det{|M_{n+1}|} = \det{|M_n|}(d - \vec{v}^\dagger M_n^{-1} \vec{v} )> 0,\]
which implies $(d - \vec{v}^\dagger M_n^{-1} \vec{v}) > 0$. 
Hence $x^\dagger M_{n+1} x > 0$ (is also positive definite). 
\qed
\end{proof}
